\section{Introduction}

This report constitutes the final assessment for the Noise and Vibrations course and addresses the vibratory behaviour of a structural geometry by means of numerical simulation. The work comprises two analysis blocks: a \textbf{modal analysis}, in which the natural frequencies and mode shapes of the selected geometry are determined, and a \textbf{harmonic analysis}, in which the structural response to a frequency-dependent sinusoidal excitation is studied, including a comparison of the effect of damping.

The methodology followed starts with the selection of a suitable geometry and its preparation in ANSYS Workbench, then proceeds through the definition of the material, the generation of a convergent mesh, the application of boundary conditions consistent with the physical problem, and the execution and interpretation of both the modal and harmonic analyses. Each modelling decision is accompanied by the corresponding technical justification.

The goal of this report is not merely the presentation of numerical results, but their \textbf{technical interpretation}: explaining what has been done, why it has been done, what assumptions have been adopted, and what the obtained results mean.